\documentclass[12pt,a4paper,fontset=fandol]{ctexart}

% 页面设置
\usepackage{geometry}
\geometry{left=2.5cm,right=2.5cm,top=2.5cm,bottom=2.5cm}

% 数学支持
\usepackage{amsmath,amssymb}

% 超链接
\usepackage{hyperref}
\hypersetup{
    colorlinks=true,
    linkcolor=blue,
    citecolor=red,
    urlcolor=blue
}

% ============================================
% GB/T 7714-2015 参考文献设置（顺序编码制）
% ============================================
\usepackage[sort&compress]{gbt7714}
\bibliographystyle{gbt7714-numerical}

% 如需切换为正文引用模式（非上标）:
% \citestyle{numbers}

% 标题信息
\title{\textbf{GB/T 7714-2015 标准参考文献示例}\\
\large （顺序编码制）}
\author{作者姓名}
\date{\today}

\begin{document}

\maketitle

\begin{abstract}
本文档演示如何在 LaTeX 中使用 \texttt{gbt7714} 宏包生成符合 GB/T 7714-2015 国家标准的中英文参考文献。
\end{abstract}

\tableofcontents
\newpage

\section{引言}

在学术写作中，参考文献的著录必须符合国家标准 GB/T 7714-2015\cite{gbt7714}。LaTeX 配合 \texttt{gbt7714} 宏包可以自动生成符合国标的参考文献列表\cite{knuth1984texbook}。

\section{理论基础}

\subsection{排版系统}

\TeX 是由 Knuth 开发的排版系统\cite{knuth1984texbook}，而 \LaTeX 是 Lamport 在其基础上开发的文档准备系统\cite{lamport1994latex}。现代 \LaTeX 发行版提供了完善的中文支持\cite{ctex2018}。

\subsection{中文支持}

中文排版需要特殊处理。CTeX 宏集提供了全面的中文支持\cite{ctex2018}，包括中文字体配置和版式调整等功能。

\section{研究方法}

本研究采用文献综述法，参考了多篇中英文文献\cite{goossens1994companion, 张三2020人工智能, 李四2018机器学习}。实验结果表明，深度学习方法在图像识别领域具有显著优势\cite{王五2019深度学习, 陈七2021自然语言}。

\section{实验结果}

如表 \ref{tab:compare} 所示，不同排版系统在处理中文文献时表现各异\cite{吴九2022知识图谱, johnson2019attention}。

\begin{table}[htbp]
\centering
\caption{排版系统对比}
\label{tab:compare}
\begin{tabular}{lcc}
\hline
\textbf{特性} & \textbf{Word} & \textbf{LaTeX} \\
\hline
公式排版 & 一般 & 优秀 \\
参考文献 & 手动 & 自动 \\
版本控制 & 困难 & 容易 \\
\hline
\end{tabular}
\end{table}

\section{结论}

LaTeX 配合 \texttt{gbt7714} 宏包是处理中文学术文献的理想工具\cite{刘十一2019博士论文, 孙十二2020硕士论文}。对于需要符合国标的学位论文和期刊投稿，建议使用 XeLaTeX 编译引擎\cite{钱十三2021技术报告}。

% 引用所有参考文献（即使未在正文中引用）
\nocite{*}

% 生成参考文献
\bibliography{references}

\end{document}
